%%%%%%%%%%%%%%%%%%%%% Chapter 26 %%%%%%%%%%%%%%%%%%%%%%%%%%%%%%%%%
% Trustworthy MLOps: Continuous Monitoring & Drift Detection
%%%%%%%%%%%%%%%%%%%%%%%%%%%%%%%%%%%%%%%%%%%%%%%%%%%%%%%%%%%%%%%%%
\motto{Trust is not a static certification; it is a dynamic performance that must be verified with every heartbeat of the system.}

\chapter{Trustworthy MLOps: Continuous Monitoring \& Drift Detection}
\label{chap:mlops_monitoring}

\abstract{A model that is verifiably trustworthy at the moment of deployment is not guaranteed to remain so as the world evolves. In this chapter, we address the operational reality of "Trustworthy MLOps"---the practice of maintaining security, privacy, and robustness in production. We explore technical mechanisms for detecting Data and Concept Drift using statistical distance metrics, analyze the unique challenges of monitoring Differential Privacy budgets in streaming environments, and define the architectural triggers for automated "Self-Healing" via rollback and retraining.}

\section{The Trust Decay Problem}
\label{sec:trust_decay}

In the "Strategic Vision" of Part I, we established that trust is built on mathematical and architectural foundations. However, even the most rigorously verified model is subject to \textbf{Performance Decay} once it encounters the entropy of the real world. 

For the Trustworthy AI Architect, "Deployment" is not the end of the project, but the beginning of a continuous monitoring cycle. Trust degrades when the statistical properties of the incoming production data shift away from the training distribution, or when the underlying relationship between inputs and outputs changes due to societal or environmental factors.

\section{Continuous Monitoring: Detecting Drift}
\label{sec:drift_detection}

To maintain trust, we must monitor two primary types of drift:

\subsection{Data Drift: Monitoring the Input}
Data drift occurs when the feature distribution of production data changes. Architects implement automated statistical tests, such as the \textbf{Population Stability Index (PSI)} or \textbf{Kullback-Leibler (K-L) Divergence}, to measure the distance between the training and production distributions ($P$ and $Q$).
\begin{equation}
D_{KL}(P \| Q) = \sum_{x \in X} P(x) \log\left(\frac{P(x)}{Q(x)}\right)
\end{equation}
If the divergence exceeds a pre-defined "Trust Threshold," the system triggers a diagnostic alert.

\subsection{Concept Drift: Monitoring the Target}
Concept drift is more insidious; it occurs when the "ground truth" changes. For example, a fraud detection model may fail as criminal tactics evolve, even if the input distribution remains the same. We use algorithms like \textbf{ADWIN (Adaptive Windowing)} to identify sudden changes in model accuracy or loss, indicating that the model's internal map of the world is no longer accurate.

\section{Privacy and Robustness Monitoring}
\label{sec:privacy_monitoring}

Beyond accuracy, we must monitor the stability of our trust pillars:

\begin{itemize}
    \item \textbf{Privacy Budget Monitoring:} For systems that use online learning or interactive queries, the architect must track the cumulative \textbf{Privacy Budget ($\epsilon$)} in real-time. If the budget approaches its limit, the system must automatically restrict query access or cease fine-tuning to prevent information leakage.
    \item \textbf{Robustness Drift:} We monitor for shifts in the success rate of our internal ``continuous red teaming'' (from Chapter 17). If the model's resilience against known adversarial patterns drops, it indicates that the model's decision boundaries have become brittle due to new data.
\end{itemize}

% ==================================================================
\section{The Privacy Audit Process}
\label{sec:privacy_audit}
% ==================================================================

Continuous monitoring (Section~\ref{sec:privacy_monitoring}) is a \emph{passive surveillance signal}; a Privacy Audit is an \emph{active certification event}. An audit produces a legally admissible, time-stamped document that answers three deterministic questions:

\begin{enumerate}
  \item \textbf{Budget Accounting:} Has the system's Differential Privacy budget $\varepsilon$ been correctly accounted and not exceeded throughout its operational lifetime?
  \item \textbf{Empirical Leakage:} Do trained models leak individual training records beyond what the $(\varepsilon, \delta)$-DP guarantee permits, as measured by black-box membership inference?
  \item \textbf{Data Flow Conformance:} Is every documented data flow consistent with the Contextual Integrity norms declared at design time (Chapter~\ref{chap:contextual_integrity}) and the purpose declarations registered under Decree~13 Art.~6?
\end{enumerate}

This audit is distinct from a general security audit or penetration test. It is specifically calibrated to produce output that satisfies \textbf{Decree 13/2023/ND-CP Article 8} (data subjects' right to processing records), \textbf{GDPR Article 5(2)} (the Accountability Principle), and \textbf{EU AI Act Article 17} (post-market monitoring obligations for High-Risk systems).

% ------------------------------------------------------------------
\subsection{Five-Phase Audit Protocol}
\label{subsec:audit_protocol}
% ------------------------------------------------------------------

The Privacy Audit is structured into five sequential phases, each with explicit entry conditions, outputs, and a responsible role. Table~\ref{tab:audit_protocol} specifies the protocol.

\begin{table}[ht]
\centering
\caption{Five-Phase Privacy Audit Protocol. The Responsible column names the architectural role accountable for each phase output. All phase outputs are inputs to the Audit Report (Phase 5).}
\label{tab:audit_protocol}
\renewcommand{\arraystretch}{1.3}
\begin{tabular}{p{0.5cm} p{2.4cm} p{4.2cm} p{2.2cm} p{2.6cm}}
\hline\noalign{\smallskip}
\textbf{\#} & \textbf{Phase} & \textbf{Activity} & \textbf{Output Artifact} & \textbf{Responsible} \\
\noalign{\smallskip}\svhline\noalign{\smallskip}
1 & Scope \& Inventory & Enumerate all data flows, models, and DP-protected components in scope; verify against the CI Norm Matrix (Ch.~\ref{chap:contextual_integrity}) & Data Flow Inventory (DFI) & Privacy Architect \\[4pt]
2 & Budget Reconciliation & Extract the Privacy Budget Ledger (Section~\ref{subsec:privacy_budget_ledger}); verify that cumulative $\varepsilon_{\text{total}} \leq \varepsilon_{\text{budget}}$ and $\delta_{\text{total}} < 1/N$ for every in-scope model & Signed Budget Report & MLOps Engineer \\[4pt]
3 & Empirical MIA Test & Execute black-box Membership Inference Attack (MIA) against each model using the protocol in Section~\ref{subsec:mia_audit}; compare observed advantage to the theoretical DP bound & MIA Advantage Score & Security Engineer \\[4pt]
4 & Data Flow Conformance & Sample 10\% of API calls from the audit period and trace each against the 5-parameter CI model; flag any flow whose Recipient or Transmission Principle deviated from the DFI & Flow Conformance Log & Data Governance \\[4pt]
5 & Report \& Filing & Synthesize phases 1--4 into the Audit Report (Section~\ref{subsec:audit_report}); file with the regulatory authority under Decree~13 Art.~8 or submit to the EU Notified Body & Signed Audit Report & DPO / Legal \\
\noalign{\smallskip}\hline\noalign{\smallskip}
\end{tabular}
\end{table}

% ------------------------------------------------------------------
\subsection{Privacy Budget Ledger}
\label{subsec:privacy_budget_ledger}
% ------------------------------------------------------------------

The \textbf{Privacy Budget Ledger} is the primary audit artifact for any system using Differential Privacy. It is an append-only, cryptographically signed log in which every mechanism invocation records a ledger entry. For a DP-FL system (Chapter~\ref{chap:fl}), one entry is created per aggregation round.

\begin{table}[ht]
\centering
\caption{Privacy Budget Ledger schema. The Cumulative column is a running sum; an automated alert fires when it crosses the pre-approved $\varepsilon_{\text{budget}}$ threshold. Each entry is signed by the TEE attestation key (Chapter~\ref{chap:tees_enclaves}) to prevent retrospective modification.}
\label{tab:budget_ledger}
\renewcommand{\arraystretch}{1.25}
\begin{tabular}{p{1.0cm} p{2.0cm} p{1.5cm} p{1.5cm} p{1.8cm} p{1.8cm} p{2.5cm}}
\hline\noalign{\smallskip}
\textbf{Round} & \textbf{Timestamp} & \textbf{$N$ clients} & \textbf{$\varepsilon_r$} & \textbf{$\delta_r$} & \textbf{$\varepsilon_{\text{cum}}$} & \textbf{TEE Signature} \\
\noalign{\smallskip}\svhline\noalign{\smallskip}
1 & 2026-03-01 08:00 & 142 & 0.12 & $10^{-6}$ & 0.12 & \texttt{3a7f...c219} \\
2 & 2026-03-01 09:15 & 137 & 0.13 & $10^{-6}$ & 0.25 & \texttt{b8e2...f04a} \\
$\vdots$ & $\vdots$ & $\vdots$ & $\vdots$ & $\vdots$ & $\vdots$ & $\vdots$ \\
$T$ & 2026-03-31 23:59 & 151 & 0.11 & $10^{-6}$ & 3.47 & \texttt{9c1d...a77b} \\
\noalign{\smallskip}\hline\noalign{\smallskip}
\end{tabular}
\end{table}

The cumulative $\varepsilon$ total is computed using the \textbf{Rényi Differential Privacy (RDP) accountant}, not simple summation, to avoid the pessimistic basic composition bound (Chapter~\ref{chap:differential_privacy}). The Ledger is implemented as an \textbf{immutable append-only log}---a blockchain or a TEE-sealed write-once store---so that neither the MLOps team nor a compromised insider can retroactively alter an $\varepsilon_r$ entry to make the budget appear unspent.

\begin{svgraybox}
\textbf{Audit Rule:} If the Privacy Budget Ledger cannot be produced and its signature chain verified, the audit \textbf{fails unconditionally}, regardless of the model's MIA test results. An unverified ledger is an unverified claim; a regulator who cannot reproduce the $\varepsilon$ computation from the raw ledger entries has no basis to certify compliance.
\end{svgraybox}

% ------------------------------------------------------------------
\subsection{Empirical Membership Inference Audit (MIA Test)}
\label{subsec:mia_audit}
% ------------------------------------------------------------------

The DP guarantee is a mathematical bound, not a measured value. An adversary with black-box access to the model may achieve a Membership Inference Attack advantage that empirically exceeds the theoretical bound if implementation errors were made (e.g., incorrect gradient clipping, batch normalization leakage, or wrong accounting mode). The empirical MIA test operationalizes the theoretical guarantee.

\subsubsection*{Protocol}

\begin{enumerate}
  \item Construct a \textbf{challenge dataset}: 1\,000 records drawn from the training set (\textbf{members}) and 1\,000 records from a held-out set with the same distribution (\textbf{non-members}).
  \item For each record $x_i$, query the model for the confidence vector $\hat{p} = f_\theta(x_i)$.
  \item Train a \textbf{shadow attack classifier} on $(x_i, \hat{p}_i, y_i)$ where $y_i \in \{\text{member}, \text{non-member}\}$.
  \item Measure the \textbf{MIA Advantage}: $\mathrm{Adv}_{\text{MIA}} = 2 \cdot (\mathrm{AUC} - 0.5)$, where AUC is the area under the ROC curve of the attack classifier.
  \item Compare against the theoretical bound. Under $(\varepsilon, \delta)$-DP, the advantage is bounded by:
\begin{equation}
  \mathrm{Adv}_{\text{MIA}} \leq e^\varepsilon - 1 + \delta \cdot e^\varepsilon
  \label{eq:mia_bound}
\end{equation}
\end{enumerate}

\begin{svgraybox}
\textbf{Pass/Fail Criterion:} If $\mathrm{Adv}_{\text{MIA}} \leq e^{\varepsilon_{\text{total}}} - 1 + \delta$, the model passes the empirical audit. If it exceeds this bound, the DP implementation contains an error and the system must be suspended pending root-cause analysis---regardless of the formal budget accounting in Phase 2.
\end{svgraybox}

% ------------------------------------------------------------------
\subsection{The Audit Report}
\label{subsec:audit_report}
% ------------------------------------------------------------------

The Audit Report is the legally significant output of the five-phase protocol. Its structure mirrors the requirements of the three governing frameworks: Decree~13 Art.~8 (audit rights), GDPR Art.~5(2) (accountability), and EU AI Act Art.~17 (post-market monitoring).

\begin{table}[ht]
\centering
\caption{Audit Report: mandatory sections and their regulatory mapping. A report missing any mandatory section is rejected by regulators under all three frameworks.}
\label{tab:audit_report}
\renewcommand{\arraystretch}{1.3}
\begin{tabular}{p{3.2cm} p{4.2cm} p{4.6cm}}
\hline\noalign{\smallskip}
\textbf{Report Section} & \textbf{Content} & \textbf{Regulatory Mapping} \\
\noalign{\smallskip}\svhline\noalign{\smallskip}
Executive Summary & System identity, audit scope, audit period, and overall Pass/Fail verdict & Required by all three frameworks as the entry point for regulatory review \\[4pt]
Budget Reconciliation & Full Privacy Budget Ledger extract; RDP accountant computation trace; final $\varepsilon_{\text{total}}$ vs.\ $\varepsilon_{\text{budget}}$; TEE signature chain & Decree 13 Art.~8 (processing records); GDPR Art.~5(2) (accountability) \\[4pt]
MIA Test Results & Challenge dataset provenance; shadow classifier AUC; $\mathrm{Adv}_{\text{MIA}}$ vs.\ theoretical bound (Eq.~\ref{eq:mia_bound}); Pass/Fail verdict & EU AI Act Art.~17 (post-market monitoring); GDPR Art.~25 (Privacy by Design) \\[4pt]
Data Flow Conformance & Flow Conformance Log summary; number of sampled calls; number of CI violations detected; violation disposition & Decree 13 Art.~6 (purpose limitation); GDPR Art.~5(1)(b) (purpose limitation) \\[4pt]
Remediation Register & Open findings, assigned owners, and deadline dates for each failed phase & EU AI Act Art.~9 (risk management system); Decree 13 Art.~26 (incident management) \\[4pt]
Signatures & DPO, Lead Auditor, and TEE-attested system timestamp & Decree 13 Art.~8; GDPR Art.~5(2) accountability obligation \\
\noalign{\smallskip}\hline\noalign{\smallskip}
\end{tabular}
\end{table}

\begin{svgraybox}
\textbf{Audit Frequency:} For \textbf{High-Risk systems} under the EU AI Act or Vietnam's AI Law, a full Privacy Audit must be conducted (1) prior to initial deployment, (2) after any model re-training or architecture change, and (3) on an annual calendar schedule. For \textbf{Medium-Risk systems}, an annual audit suffices. For \textbf{Low-Risk systems}, a lightweight budget reconciliation (Phase 2 only) is sufficient on a semi-annual basis.
\end{svgraybox}

% ------------------------------------------------------------------
%  Colour definitions for Trust Debt section
% ------------------------------------------------------------------
% Colour definitions consolidated in book.tex preamble

% ------------------------------------------------------------------
%  tcolorbox styles for Trust Debt section
% ------------------------------------------------------------------
\tcbset{
  scenariobox/.style={
    enhanced, colback=scenariobg,
    colframe=gray!60, fonttitle=\bfseries\small,
    left=8pt, right=8pt, top=6pt, bottom=6pt,
    breakable
  },
  warningbox/.style={
    enhanced, colback=warningbg,
    colframe=debtamber!80!black, fonttitle=\bfseries\small,
    left=8pt, right=8pt, top=6pt, bottom=6pt,
    breakable
  },
  debtpattern/.style={
    enhanced, colback=blue!4,
    colframe=blue!55!black, fonttitle=\bfseries\small,
    title={Debt Management Pattern},
    left=8pt, right=8pt, top=6pt, bottom=6pt,
    breakable
  }
}

% ==================================================================
\section{Managing Trustworthy AI Technical Debt}
\label{sec:trust-tech-debt}
% ==================================================================

\begin{quote}
\textit{``Complexity is the enemy of security. Every layer of
trust you add to a system is also a layer of debt that must be
serviced, versioned, and eventually repaid. The question is not
whether your ZK-circuit will need an update---it is whether you
will be ready when it does.''}
\end{quote}

Sections~\ref{sec:drift_detection}
and~\ref{sec:privacy_monitoring} established the
monitoring infrastructure for detecting drift in model accuracy,
privacy budgets, and adversarial resilience.
What those sections did not address is the deeper operational
question that monitoring eventually surfaces: when the
monitoring system tells you that a trust mechanism is failing,
degraded, or exhausted, \emph{how do you fix it without breaking
everything else?}

This is the problem of \textbf{Trustworthy AI Technical Debt}:
the accumulated cost of maintaining, updating, and replacing
interconnected trust mechanisms in a live production system,
where each mechanism was designed with specific assumptions
about the others, and where a change to any single component
propagates consequences through the entire stack.

Standard software technical debt is well-understood.
Trustworthy AI technical debt has three properties that make it
substantially more dangerous:

\begin{description}
  \item[Cryptographic Coupling.]
    Trust mechanisms are not loosely coupled services.
    A ZK-circuit compiled from a specific model checkpoint,
    a TEE enclave signed against a specific firmware version,
    and a DP privacy budget calibrated to a specific dataset
    size are all \emph{mathematically bound} to each other.
    Updating one without updating the others invalidates
    the system's formal guarantees---silently, with no
    runtime error.

  \item[Invisible Degradation.]
    Unlike a crashed microservice, a degraded trust mechanism
    continues to function.
    A ZK-circuit running against a model that has been
    fine-tuned since the circuit was compiled still produces
    proofs---but those proofs are now proofs of the
    \emph{wrong computation}.
    The system appears healthy while its trustworthiness
    guarantee has been severed.

  \item[Regulatory Irreversibility.]
    Under Vietnam's AI Law and the EU AI Act, a High-Risk
    system that has been operating with an invalidated trust
    guarantee may be required to suspend operations, notify
    regulators, and resubmit for Conformity Assessment---even
    if the invalidity was unintentional and caused no
    measurable harm.
    Technical debt in a Trustworthy AI system is therefore
    not just an engineering liability; it is a regulatory
    exposure.
\end{description}

This section provides a formal taxonomy of Trustworthy AI
technical debt, four management patterns for the most critical
debt scenarios, and a Debt Registry implementation that gives
the architect continuous visibility into the health of every
trust mechanism in the stack.

% ------------------------------------------------------------------
\subsection{The Trustworthy AI Debt Taxonomy}
\label{subsec:debt-taxonomy}
% ------------------------------------------------------------------

We define \textbf{Trustworthy AI Technical Debt} as any
condition in which the deployed trust mechanism stack no longer
provides the formal guarantees that were certified at the time
of the last Conformity Assessment.
Debt arises from four distinct sources, each with a distinct
accumulation profile, detection method, and repayment strategy.

\begin{description}
  \item[Type~1: Budget Debt.]
    A consumable resource---most commonly the Differential
    Privacy budget $\varepsilon$---has been partially or fully
    exhausted through normal operation.
    Budget debt accumulates continuously and predictably.
    It is the only debt type that can be precisely tracked in
    real time without additional tooling.

  \item[Type~2: Staleness Debt.]
    A trust mechanism has been compiled, trained, or certified
    against an artifact that has since changed.
    The most common instances are a ZK-circuit compiled against
    a model checkpoint that has been fine-tuned, and a Formal
    Verification proof computed against a model architecture
    that has since been pruned or quantised.
    Staleness debt is discrete: the mechanism is either current
    or stale; there is no intermediate state.

  \item[Type~3: Composition Debt.]
    The formal guarantees of a multi-mechanism stack depend
    on specific assumptions about the interaction between
    mechanisms.
    Composition debt arises when one mechanism in the stack
    is updated in a way that invalidates the composition
    proof of the overall stack, even though each individual
    mechanism remains internally valid.
    This is the most dangerous debt type because it is
    invisible to mechanism-level monitoring.

  \item[Type~4: Entropy Debt.]
    Cryptographic mechanisms have finite operational lifetimes
    determined by advances in computational power and
    cryptanalysis.
    A TEE enclave signed with a firmware key that has since
    been disclosed, or a ZK-proof system whose underlying
    elliptic curve has been weakened by a new attack, carries
    entropy debt.
    Entropy debt accumulates slowly but arrives catastrophically:
    the mechanism is secure until the day it is not.
\end{description}

Table~\ref{tab:debt-taxonomy} summarises the four debt types
with their accumulation profiles, primary detection methods,
and time-to-criticality estimates for a typical Vietnamese
enterprise AI deployment.

\begin{table}[ht]
\centering
\caption{Trustworthy AI Technical Debt Taxonomy.
  Time-to-criticality is the estimated elapsed time from
  debt inception to regulatory or security failure in the
  absence of active management, based on observed production
  patterns and the Vietnamese AI Law Conformity Assessment
  cycle.}
\label{tab:debt-taxonomy}
\small
\begin{tabular}{p{1.8cm}p{2.2cm}p{2.0cm}p{2.8cm}p{2.2cm}}
\hline\noalign{\smallskip}
\textbf{Debt Type}
  & \textbf{Accumulation Profile}
  & \textbf{Visibility}
  & \textbf{Primary Detection Method}
  & \textbf{Time-to-Criticality} \\
\noalign{\smallskip}\svhline\noalign{\smallskip}

Type~1 Budget
  & Continuous, predictable
  & \cellcolor{debtgreen!30} High
  & Real-time $\varepsilon$ tracking (Section~\ref{sec:privacy_monitoring})
  & Weeks to months \\[6pt]

Type~2 Staleness
  & Discrete, event-triggered
  & \cellcolor{debtamber!40} Medium
  & Artifact hash comparison in Debt Registry
  & Hours to days after trigger \\[6pt]

Type~3 Composition
  & Discrete, hidden
  & \cellcolor{debtred!30} Low
  & Stack-level composition audit (quarterly)
  & Days to weeks after trigger \\[6pt]

Type~4 Entropy
  & Slow, then sudden
  & \cellcolor{debtred!30} Low
  & CVE monitoring + hardware vendor advisories
  & Months to years \\

\noalign{\smallskip}\hline\noalign{\smallskip}
\end{tabular}
\end{table}

% ------------------------------------------------------------------
\subsection{The Trust Debt Accumulation Model}
\label{subsec:debt-accumulation}
% ------------------------------------------------------------------

We model the aggregate \textbf{Trust Debt Score} $\mathcal{D}(t)$
of a deployed system at time $t$ as a weighted sum of
active debt indicators across all four debt types and all
mechanisms in the stack $M = \{m_1, \ldots, m_k\}$:

\begin{equation}
  \label{eq:debt-score}
  \mathcal{D}(t)
    = \sum_{j=1}^{4} w_j
      \sum_{i=1}^{k}
        \mathbf{1}\!\left[
          \delta_{j}(m_i, t) > \theta_j
        \right]
      \cdot s_{j}(m_i, t)
\end{equation}

where:
\begin{itemize}
  \item $w_j \in [0,1]$ is the regulatory weight of debt
    type $j$, with $\sum_j w_j = 1$.
    Recommended weights for a Vietnamese High-Risk deployment:
    $w_1 = 0.30$ (Budget),
    $w_2 = 0.25$ (Staleness),
    $w_3 = 0.35$ (Composition),
    $w_4 = 0.10$ (Entropy).

  \item $\delta_{j}(m_i, t)$ is the debt indicator for
    mechanism $m_i$ of type $j$ at time $t$
    (defined per-type in the scenarios below).

  \item $\theta_j$ is the alert threshold for debt type $j$
    above which the indicator contributes to the score.

  \item $s_{j}(m_i, t) \in [0, 1]$ is the severity of the
    debt instance, normalised to $[0,1]$.
\end{itemize}

A Debt Score $\mathcal{D}(t) = 0$ indicates a fully current
stack with no active debt.
$\mathcal{D}(t) \geq 1.0$ indicates a stack-level failure
requiring immediate operational suspension under the Emergency
Brake Pattern of Section~\ref{sec:agency_spectrum}.

Figure~\ref{fig:debt-accumulation} illustrates a representative
annual debt accumulation profile showing all four debt types
and the system's response events.

\begin{figure}[ht]
\centering
\begin{tikzpicture}
\begin{axis}[
    width=12cm, height=6.5cm,
    xlabel={Month of deployment year},
    ylabel={Trust Debt Score $\mathcal{D}(t)$},
    xmin=0, xmax=12,
    ymin=0, ymax=1.35,
    xtick={0,1,2,3,4,5,6,7,8,9,10,11,12},
    xticklabels={Jan,Feb,Mar,Apr,May,Jun,Jul,Aug,Sep,Oct,Nov,Dec,},
    ytick={0, 0.25, 0.50, 0.75, 1.00, 1.25},
    grid=major, grid style={dashed, gray!35},
    tick label style={font=\footnotesize},
    label style={font=\small},
    legend style={
      at={(0.02,0.98)}, anchor=north west,
      font=\footnotesize, fill=white, draw=gray!50
    },
  ]

  %% Emergency suspension threshold
  \addplot[debtred, dashed, thick, domain=0:12] {1.0}
      node[right, font=\footnotesize, debtred] {Emergency Brake};

  %% Alert threshold
  \addplot[debtamber, dashed, domain=0:12] {0.60}
      node[right, font=\footnotesize, debtamber] {Alert threshold};

  %% Composite debt score
  \addplot[blue!70!black, thick, mark=none] coordinates {
    (0,  0.00)
    (1,  0.08)
    (2,  0.17)
    (3,  0.30)
    (3,  0.55)
    (4,  0.62)
    (4.5,0.35)
    (5,  0.28)
    (6,  0.38)
    (7,  0.60)
    (7,  0.85)
    (7.5,0.45)
    (8,  0.38)
    (9,  0.48)
    (10, 0.55)
    (10.5,0.15)
    (11, 0.20)
    (12, 0.28)
  };
  \addlegendentry{Composite $\mathcal{D}(t)$}

  %% Remediation event markers
  \addplot[mark=triangle*, mark size=3pt, only marks,
           debtgreen, mark options={fill=debtgreen}] coordinates {
    (4.5, 0.35)
    (7.5, 0.45)
    (10.5,0.15)
  };
  \addlegendentry{Remediation event}

  %% Debt spike markers
  \addplot[mark=*, mark size=2.5pt, only marks,
           debtred, mark options={fill=debtred}] coordinates {
    (3,  0.55)
    (7,  0.85)
  };
  \addlegendentry{Debt spike (trigger event)}

  %% Annotation: staleness event
  \node[font=\scriptsize, align=center, debtred]
      at (axis cs: 3.2, 0.65)
      {Model\\fine-tuned};

  %% Annotation: composition debt
  \node[font=\scriptsize, align=center, debtred]
      at (axis cs: 7.5, 0.96)
      {Composition\\debt detected};

  %% Annotation: budget reset
  \node[font=\scriptsize, align=center, debtgreen]
      at (axis cs: 10.8, 0.28)
      {Budget\\reset};

\end{axis}
\end{tikzpicture}
\caption{Representative annual Trust Debt Score accumulation
  profile for a Vietnamese High-Risk AI deployment.
  Budget debt (Type~1) accumulates continuously between
  remediation events.
  The model fine-tuning event at Month~3 triggers a sudden
  staleness debt (Type~2) spike.
  The composition debt event at Month~7 (Type~3) pushes the
  system into the alert zone, requiring an emergency stack
  revalidation.
  A quarterly DP budget reset at Month~10 returns the score
  to its lowest point of the year.}
\label{fig:debt-accumulation}
\end{figure}

% ------------------------------------------------------------------
\subsection{Four Critical Debt Scenarios and Their Management
            Patterns}
\label{subsec:debt-scenarios}
% ------------------------------------------------------------------

\subsubsection{Scenario 1: The Exhausted DP Budget}
\label{subsubsec:scenario-dp-budget}

\begin{tcolorbox}[scenariobox,
  title={Scenario 1 --- The Exhausted DP Budget (Type~1 Debt)}]

\textbf{Trigger:}
The cumulative privacy loss tracker
(Section~\ref{sec:privacy_monitoring})
reports that the quarterly DP budget $\varepsilon_{\mathrm{consumed}}$
has reached or exceeded the certified budget
$\varepsilon_{\mathrm{budget}}$:
\begin{equation}
  \label{eq:budget-exhaustion}
  \varepsilon_{\mathrm{consumed}}(t)
    = \sum_{\tau=0}^{t} \varepsilon_{\tau}
    \geq \varepsilon_{\mathrm{budget}}
\end{equation}

\textbf{Consequence if unmanaged:}
Any further queries or fine-tuning steps consume privacy budget
beyond the certified level.
The system continues to operate but its $(\varepsilon,
\delta)$-DP guarantee is now invalid.
Under Decree~13, this constitutes a privacy control failure
that must be reported within 72 hours of detection.

\textbf{Why this happens faster than expected:}
Three operational patterns accelerate budget consumption beyond
design estimates: (1)~unplanned A/B testing against live data;
(2)~emergency fine-tuning triggered by accuracy drift alerts
(Section~\ref{sec:drift_detection}); (3)~interactive query
APIs that were not included in the original budget allocation.
Each of these is a legitimate operational action that, in
aggregate, exhausts the budget ahead of the quarterly reset.
\end{tcolorbox}

\begin{tcolorbox}[debtpattern,
  title={Debt Management Pattern 1 --- Budget Triage Protocol}]

The Budget Triage Protocol defines a three-stage response to
approaching budget exhaustion, triggered at three threshold
levels:

\medskip
\begin{tabular}{@{}p{1.6cm} p{2.0cm} p{7.2cm}@{}}
\hline\noalign{\smallskip}
\textbf{Stage} & \textbf{Threshold} & \textbf{Actions} \\
\noalign{\smallskip}\svhline\noalign{\smallskip}
\textsc{Yellow}
  & $\varepsilon_{\mathrm{consumed}} \geq 0.70\,
    \varepsilon_{\mathrm{budget}}$
  & Disable all non-essential interactive queries.
    Freeze scheduled fine-tuning.
    Notify Data Protection Officer.
    Project exhaustion date at current burn rate. \\[4pt]
\textsc{Orange}
  & $\varepsilon_{\mathrm{consumed}} \geq 0.90\,
    \varepsilon_{\mathrm{budget}}$
  & Switch all inference to read-only mode
    (no gradient computation).
    Activate cached response serving for common queries.
    DPO escalation: request emergency budget reallocation
    or regulatory guidance. \\[4pt]
\textsc{Red}
  & $\varepsilon_{\mathrm{consumed}} \geq
    \varepsilon_{\mathrm{budget}}$
  & Suspend all model interactions that consume budget.
    Serve only cached, pre-certified responses.
    File regulatory notification within 72 hours.
    Initiate emergency budget reset procedure. \\
\noalign{\smallskip}\hline\noalign{\smallskip}
\end{tabular}

\medskip
\textbf{Budget Reset Procedure:}
A DP budget reset is not simply incrementing a counter.
It requires: (1)~generating a new Privacy Accountant checkpoint
with a fresh composition baseline; (2)~revalidating the
Gaussian noise scale $\sigma$ against the new
$(\varepsilon_{\mathrm{new}}, \delta)$ target, accounting for
any model or dataset changes since the last calibration;
(3)~updating the Conformity Assessment documentation to record
the reset event, its trigger, and the actions taken during
each triage stage.
\end{tcolorbox}

\subsubsection{Scenario 2: The Stale ZK-Circuit}
\label{subsubsec:scenario-zk-circuit}

\begin{tcolorbox}[scenariobox,
  title={Scenario 2 --- The Stale ZK-Circuit (Type~2 Debt)}]

\textbf{Trigger:}
The production model has been updated---through fine-tuning,
quantisation, or weight pruning---since the ZK-circuit was
compiled.
The ZK-proof system continues to generate valid proofs, but
those proofs now verify execution of the \emph{old} model
computation, not the current one.

\textbf{Formal statement of the failure:}
Let $\mathcal{C}_{\mathrm{circuit}}$ be the arithmetic circuit
compiled from model checkpoint $\theta_0$, and let
$\theta_1 \neq \theta_0$ be the current production weights.
The ZK-proof $\pi$ generated for a query $(x, y)$ satisfies:
\begin{equation}
  \label{eq:stale-circuit}
  \mathrm{Verify}(\pi, \mathcal{C}_{\mathrm{circuit}}, x, y)
    = \textsc{True}
  \quad \text{but} \quad
  y \neq f_{\theta_1}(x)
\end{equation}
The verifier accepts the proof, but the proof is a proof of
the wrong function.
The system has \emph{verifiable incorrectness}.

\textbf{Why this is common:}
Model updates in production are frequent (accuracy drift
remediation, fairness corrections, security patches) and are
typically handled by the MLOps pipeline with no awareness of
the ZK-circuit compilation dependency.
The circuit compilation step is expensive (hours of GPU time
for a 7B-parameter model) and is therefore omitted from
the fast-path update workflow.
This is the canonical form of Trustworthy AI technical debt:
the expensive trust component is the first to be skipped
under time pressure.
\end{tcolorbox}

\begin{tcolorbox}[debtpattern,
  title={Debt Management Pattern 2 --- Circuit Versioning Gate}]

The Circuit Versioning Gate is a mandatory CI/CD pipeline
check that blocks any model weight update from reaching
production unless one of two conditions is satisfied:
(1)~the update is certified as \emph{circuit-preserving} by
a structural equivalence check; or (2)~a new ZK-circuit has
been compiled and validated against the updated weights.

\textbf{Circuit-Preserving Update Definition:}
An update $\theta_0 \to \theta_1$ is circuit-preserving if
and only if the model's architectural graph---the sequence
of operations, their dimensions, and their activation
functions---is unchanged.
Weight-only updates (e.g., gradient descent steps, LoRA
adapter updates where the base architecture is frozen)
are circuit-preserving.
Structural changes (quantisation bit-width changes, layer
pruning, architectural modifications) are circuit-breaking
and require circuit recompilation.

\textbf{Implementation:}
The gate is implemented as a Git pre-merge hook that computes
the SHA-256 hash of the model's ONNX computational graph
(structure only, not weights) and compares it to the hash
recorded in the Debt Registry at the time of the last
circuit compilation.
A hash mismatch blocks the merge and opens a Debt Registry
ticket of Type~2 with estimated recompilation cost.
\end{tcolorbox}

\subsubsection{Scenario 3: The Broken Composition Guarantee}
\label{subsubsec:scenario-composition}

\begin{tcolorbox}[scenariobox,
  title={Scenario 3 --- The Broken Composition Guarantee
         (Type~3 Debt)}]

\textbf{Trigger:}
A component of the trust stack has been independently updated
in a way that violates the assumptions of the composition
proof.
The canonical example: the Secure Aggregation protocol is
upgraded to a new batched variant, changing the masking
granularity.
The existing DP noise calibration was computed under the
assumption of the old SecAgg's aggregation properties.
The new SecAgg has different sensitivity bounds.
The DP guarantee $(\varepsilon, \delta)$ that the stack
certified is no longer achievable at the configured noise
level $\sigma$.

\textbf{The composition theorem dependency graph:}
The formal guarantee of the full stack
(TEE + DP + SecAgg) is not the conjunction
of three independent guarantees.
It is a \emph{proof over their interaction}, which makes
specific assumptions about each component's operational
parameters.
Updating any component breaks the proof until it is
re-derived for the new configuration.

\textbf{Why this is the most dangerous debt type:}
Unlike Scenario~1 (budget exhaustion is visible) and
Scenario~2 (circuit hash mismatch is detectable),
composition debt produces no observable signal in
mechanism-level monitoring.
Every individual component passes its own health checks.
The stack as a whole is broken.
\end{tcolorbox}

\begin{tcolorbox}[debtpattern,
  title={Debt Management Pattern 3 --- Composition Audit Cycle}]

The Composition Audit Cycle is a mandatory quarterly review
process that re-derives the stack-level composition guarantee
from the current operational parameters of each mechanism.

\textbf{Audit Inputs:}
\begin{itemize}
  \item Current DP noise scale $\sigma$ and clipping norm $C$
    for each DP-protected component.
  \item Current SecAgg aggregation group size $n$ and
    masking scheme version.
  \item Current TEE firmware version and attestation
    certificate expiry.
  \item Current ZK-circuit version and proof system
    parameters (elliptic curve, field size).
\end{itemize}

\textbf{Audit Procedure:}
The audit re-derives the effective privacy guarantee
$\varepsilon_{\mathrm{eff}}$ of the composed stack using
the Privacy Loss Random Variable (PRV) accountant or
R\'{e}nyi DP composition, and
compares it to the certified value $\varepsilon_{\mathrm{cert}}$.
If $\varepsilon_{\mathrm{eff}} > \varepsilon_{\mathrm{cert}}$,
the stack carries active Type~3 debt that must be remediated
before the next production traffic window.

\textbf{Audit Trigger Events:}
In addition to the quarterly cycle, a Composition Audit is
triggered immediately by any of the following events:
(1)~SecAgg protocol version change;
(2)~TEE firmware update;
(3)~DP accountant library version update;
(4)~addition or removal of any mechanism from the stack;
(5)~any change to the production dataset size $N$ that
affects the $\delta = 1/N$ calibration.

\textbf{Composition Debt Resolution:}
Resolution requires recalibrating $\sigma$ to restore
$\varepsilon_{\mathrm{eff}} \leq \varepsilon_{\mathrm{cert}}$
under the new composition, or obtaining a new Conformity
Assessment certification for the updated $\varepsilon_{\mathrm{cert}}$.
The recalibrated stack must pass the full acceptance
protocol before returning to production.
\end{tcolorbox}

\subsubsection{Scenario 4: The Decaying TEE Enclave}
\label{subsubsec:scenario-tee-entropy}

\begin{tcolorbox}[scenariobox,
  title={Scenario 4 --- The Decaying TEE Enclave (Type~4 Debt)}]

\textbf{Trigger:}
One of three entropy events: (1)~Intel or AMD publishes a
microcode security advisory (SGX~AEPIC leak class,
Spectre/Meltdown variants) requiring enclave firmware
patching; (2)~the attestation certificate's root of trust
key approaches expiry; (3)~a side-channel attack technique
is published that targets the specific memory access pattern
of the deployed model workload (as occurred in the LedgerNet
incident, Appendix~\ref{app:roi-of-trust}).

\textbf{The operational complexity of TEE patching:}
A TEE firmware patch is not a software update.
It requires: (1)~draining all live inference requests from
the enclave; (2)~re-measuring the enclave image against
the new firmware; (3)~re-generating and distributing the
remote attestation certificate to all clients that hold
the previous certificate; (4)~re-establishing all encrypted
channels that were bound to the old attestation.
For a system serving thousands of concurrent sessions,
this is a multi-hour operation with complete service
interruption.

\textbf{The compounding risk:}
Entropy debt does not announce itself with a monitoring
alert.
It arrives via an external CVE or vendor advisory and
requires an immediate response to a complex, high-stakes
operation on a live production system---precisely the
conditions under which human error rates are highest and
the consequences of error are most severe.
\end{tcolorbox}

\begin{tcolorbox}[debtpattern,
  title={Debt Management Pattern 4 --- Enclave Lifecycle
         Management Protocol}]

The Enclave Lifecycle Management Protocol (ELMP) pre-plans
the full TEE update workflow so that it can be executed
under time pressure without improvisation.

\textbf{Standing Maintenance Window:}
Every deployment that uses TEEs must have a monthly
maintenance window of at least 4 hours during which the
enclave can be patched without violating SLA commitments.
The window schedule and the fallback serving strategy (cached
responses, graceful degradation to non-enclave inference
with appropriate user notification) must be documented in
the system's operational runbook.

\textbf{CVE Response Ladder:}
\begin{enumerate}
  \item \textbf{CVE Published (Day~0):}
    Security team assesses applicability to deployed
    enclave configuration within 4 hours.
    If applicable, the maintenance window is activated
    immediately regardless of schedule.

  \item \textbf{Patch Staging (Day~0--1):}
    New enclave image built and measured in staging
    environment.
    New attestation certificate generated and verified
    against staging clients.
    Composition Audit triggered (Pattern~3) to confirm
    that the patched enclave does not invalidate stack
    composition guarantees.

  \item \textbf{Canary Deployment (Day~1--2):}
    Patched enclave deployed to 5\% of production traffic.
    Attestation success rate, inference latency, and
    ZK-proof validity monitored for 24 hours.

  \item \textbf{Full Rollout (Day~2--3):}
    Patched enclave promoted to 100\% of production traffic.
    Old enclave drained and decommissioned.
    Debt Registry updated: Type~4 debt entry closed.
    Conformity Assessment documentation updated with
    patch event record.
\end{enumerate}

\textbf{Certificate Rotation:}
Remote attestation certificates must be rotated before
expiry on a schedule maintained in the Debt Registry.
The recommended rotation trigger is 30 days before
certificate expiry, providing a buffer for the rotation
workflow to complete within the standing maintenance window.
\end{tcolorbox}

% ------------------------------------------------------------------
\subsection{The Trust Debt Registry}
\label{sec:trust-debt-registry}
% ------------------------------------------------------------------

The four debt management patterns of
Section~\ref{subsec:debt-scenarios} share a common
operational requirement: a single authoritative record of
every active and resolved debt item across the full trust
stack.
We call this the \textbf{Trust Debt Registry}: a versioned,
auditable data structure maintained as a first-class
infrastructure artifact alongside the model weights, the
ZK-circuit, and the DP accountant state.

The Trust Debt Registry serves three functions:

\begin{enumerate}
  \item \textbf{Continuous Debt Scoring.}
    Provides the input data for the Trust Debt Score
    $\mathcal{D}(t)$ of
    Equation~\eqref{eq:debt-score}, enabling real-time
    dashboard display.

  \item \textbf{Regulatory Evidence.}
    Provides a tamper-evident record of every debt event,
    its detection time, the actions taken, and the
    resolution timestamp---the primary evidence artifact
    for Conformity Assessment re-submissions and regulatory
    audits under Vietnam's AI Law.

  \item \textbf{Remediation Scheduling.}
    Drives the operational calendar: circuit recompilation
    windows, composition audit dates, enclave maintenance
    windows, and budget reset checkpoints are all derived
    from open Debt Registry entries.
\end{enumerate}

Listing~\ref{lst:debt-registry} provides a complete
implementation of the Trust Debt Registry, including the
Debt Score calculator, the four debt detectors, and the
registry persistence layer.

\begin{lstlisting}[language=Python, style=pythonstyle,
  caption={Trust Debt Registry: unified tracking and scoring
           for all four Trustworthy AI debt types across the
           full mechanism stack.},
  label={lst:debt-registry}
]
from dataclasses import dataclass, field
from enum import Enum
from typing import Optional
import hashlib, time, json

# -- Debt taxonomy ------------------------------------------------
class DebtType(Enum):
    BUDGET      = 1   # Consumable resource exhaustion
    STALENESS   = 2   # Artifact hash mismatch
    COMPOSITION = 3   # Stack-level guarantee invalidation
    ENTROPY     = 4   # Cryptographic / hardware decay

class DebtStatus(Enum):
    OPEN     = "open"
    TRIAGED  = "triaged"
    RESOLVED = "resolved"

@dataclass
class DebtEntry:
    debt_id:      str
    debt_type:    DebtType
    mechanism:    str            # e.g., "DP-SGD", "ZK-circuit-v3"
    description:  str
    severity:     float          # normalised [0, 1]
    detected_at:  float = field(default_factory=time.time)
    resolved_at:  Optional[float] = None
    status:       DebtStatus = DebtStatus.OPEN
    audit_hash:   Optional[str] = None   # ZK-verified audit ref

# -- Registry weights for High-Risk Vietnamese deployment ---------
DEBT_WEIGHTS = {
    DebtType.BUDGET:      0.30,
    DebtType.STALENESS:   0.25,
    DebtType.COMPOSITION: 0.35,
    DebtType.ENTROPY:     0.10,
}

class TrustDebtRegistry:
    def __init__(self, audit_logger, alert_threshold=0.60,
                 emergency_threshold=1.00):
        self.entries   = []
        self.auditor   = audit_logger
        self.alert_thr = alert_threshold
        self.emerg_thr = emergency_threshold

    # -- Core CRUD ------------------------------------------------
    def open_debt(self, debt_type, mechanism,
                  description, severity):
        entry = DebtEntry(
            debt_id     = self._generate_id(mechanism, debt_type),
            debt_type   = debt_type,
            mechanism   = mechanism,
            description = description,
            severity    = min(max(severity, 0.0), 1.0),
        )
        entry.audit_hash = self.auditor.log(
            event="debt_opened", entry=entry.__dict__)
        self.entries.append(entry)
        return entry

    def resolve_debt(self, debt_id, resolution_note):
        for entry in self.entries:
            if entry.debt_id == debt_id \
                    and entry.status != DebtStatus.RESOLVED:
                entry.status      = DebtStatus.RESOLVED
                entry.resolved_at = time.time()
                self.auditor.log(
                    event="debt_resolved",
                    debt_id=debt_id,
                    note=resolution_note)
                return
        raise ValueError(f"Debt {debt_id} not found/resolved.")

    # -- Debt Score (Equation 25.1) --------------------------------
    def compute_score(self):
        open_entries = [e for e in self.entries
                        if e.status != DebtStatus.RESOLVED]
        score = sum(
            DEBT_WEIGHTS[e.debt_type] * e.severity
            for e in open_entries
        )
        return round(score, 4)

    # -- Automated Debt Detectors ---------------------------------
    def check_budget_debt(self, eps_consumed, eps_budget,
                          mechanism="DP-SGD"):
        """Type 1: Open debt if budget >= 70%."""
        ratio = eps_consumed / eps_budget
        if ratio >= 0.70:
            severity = min((ratio - 0.70) / 0.30, 1.0)
            self.open_debt(
                DebtType.BUDGET, mechanism,
                f"Budget at {ratio:.0%}: "
                f"eps={eps_consumed:.3f}/{eps_budget:.3f}",
                severity=severity)

    def check_staleness_debt(self, current_onnx_graph,
                             certified_hash,
                             mechanism="ZK-circuit"):
        """Type 2: Open debt if model graph hash changed."""
        current_hash = hashlib.sha256(
            current_onnx_graph).hexdigest()
        if current_hash != certified_hash:
            self.open_debt(
                DebtType.STALENESS, mechanism,
                f"ONNX hash mismatch: cert={certified_hash[:12]}"
                f"... cur={current_hash[:12]}...",
                severity=0.85)

    def check_composition_debt(self, eps_effective,
                               eps_certified,
                               mechanism="DP+SecAgg"):
        """Type 3: Open debt if eps_eff > eps_cert."""
        if eps_effective > eps_certified:
            severity = min(
                (eps_effective - eps_certified) / eps_certified,
                1.0)
            self.open_debt(
                DebtType.COMPOSITION, mechanism,
                f"Composition violated: "
                f"eps_eff={eps_effective:.3f} > "
                f"eps_cert={eps_certified:.3f}",
                severity=severity)

    def check_entropy_debt(self, cert_expiry_days,
                           rotation_buffer_days=30,
                           mechanism="TEE-attestation"):
        """Type 4: Open debt if certificate nears expiry."""
        if cert_expiry_days <= rotation_buffer_days:
            severity = 1.0 - (cert_expiry_days /
                               rotation_buffer_days)
            self.open_debt(
                DebtType.ENTROPY, mechanism,
                f"Cert expires in {cert_expiry_days} days",
                severity=severity)

    # -- Operational Interface ------------------------------------
    def operational_status(self):
        score  = self.compute_score()
        status = ("NOMINAL"   if score < self.alert_thr  else
                  "ALERT"     if score < self.emerg_thr  else
                  "EMERGENCY")
        return {
            "debt_score":   score,
            "status":       status,
            "open_entries": sum(1 for e in self.entries
                                if e.status != DebtStatus.RESOLVED),
            "by_type": {
                dt.name: sum(1 for e in self.entries
                             if e.debt_type == dt
                             and e.status != DebtStatus.RESOLVED)
                for dt in DebtType
            }
        }

    @staticmethod
    def _generate_id(mechanism, debt_type):
        raw = f"{mechanism}-{debt_type.name}-{time.time()}"
        return "DEBT-" + hashlib.sha256(
            raw.encode()).hexdigest()[:8].upper()
\end{lstlisting}

% ------------------------------------------------------------------
\subsection{The Debt-Aware Deployment Checklist}
\label{subsec:deployment-checklist}
% ------------------------------------------------------------------

The four debt management patterns and the Debt Registry
implementation converge in a single operational artifact:
the \textbf{Debt-Aware Deployment Checklist}, which must
be completed and signed off by both the Technical Lead and
the Data Protection Officer before any model or mechanism
update is promoted to production.

\begin{table}[ht]
\centering
\caption{Debt-Aware Deployment Checklist.
  All items must be \textbf{PASS} before the deployment
  proceeds.
  Any \textbf{FAIL} or \textbf{BLOCK} opens a mandatory
  Debt Registry entry and suspends the deployment until
  the entry is resolved.}
\label{tab:deployment-checklist}
\small
\begin{tabular}{p{0.4cm}p{4.4cm}p{2.0cm}p{4.0cm}}
\hline\noalign{\smallskip}
\textbf{\#}
  & \textbf{Check}
  & \textbf{Debt Type}
  & \textbf{Pass Condition} \\
\noalign{\smallskip}\svhline\noalign{\smallskip}

1 & DP budget remaining $\geq 30\%$ of quarterly allocation
  & Type~1
  & $\varepsilon_{\mathrm{remaining}} \geq 0.30\,
    \varepsilon_{\mathrm{budget}}$ \\[3pt]

2 & ONNX graph hash matches certified circuit hash
  & Type~2
  & SHA-256 match confirmed \\[3pt]

3 & LoRA / fine-tune delta is circuit-preserving
  & Type~2
  & Architectural graph unchanged \\[3pt]

4 & Composition audit current ($\leq 90$ days old)
  & Type~3
  & Audit timestamp in Registry \\[3pt]

5 & $\varepsilon_{\mathrm{eff}} \leq \varepsilon_{\mathrm{cert}}$
    under new configuration
  & Type~3
  & PRV accountant output verified \\[3pt]

6 & No open CVEs affecting TEE firmware version
  & Type~4
  & CVE scan result: clean \\[3pt]

7 & Attestation certificate valid $> 30$ days
  & Type~4
  & Expiry date confirmed in Registry \\[3pt]

8 & Debt Score $\mathcal{D}(t) < 0.60$
  & All types
  & Registry operational status: NOMINAL \\[3pt]

9 & ZK-circuit recompiled if architectural change detected
  & Type~2
  & New circuit hash recorded in Registry \\[3pt]

10 & DPO sign-off recorded in Registry
   & Regulatory
   & DPO signature timestamp present \\

\noalign{\smallskip}\hline\noalign{\smallskip}
\end{tabular}
\end{table}

% ------------------------------------------------------------------
\subsection{Conclusion: Debt as a First-Class Architectural
            Concern}
\label{subsec:debt-conclusion}
% ------------------------------------------------------------------

The four debt types of Section~\ref{subsec:debt-taxonomy}
represent a fundamental property of the Trustworthy AI Stack:
it is not a static certificate but a dynamic commitment that
must be continuously earned.
Every mechanism in the stack---every DP noise calibration,
every ZK-circuit, every TEE attestation, every composition
proof---carries an implicit expiry condition.
The architect who does not track these conditions is not
building a Trustworthy AI system.
They are building a system that was trustworthy on the day
of its Conformity Assessment and has been accruing hidden
debt ever since.

The Trust Debt Registry, the four management patterns, and
the Deployment Checklist of this section are not
bureaucratic additions to the engineering workflow.
They are the operational expression of the most important
insight of this book: \textbf{trust is a performance, not a
property.}
It must be demonstrated, measured, and maintained with every
deployment, every update, and every operational decision.

The Self-Healing Lifecycle of Section~\ref{sec:self_healing}
provides the automated remediation responses for the subset
of debt conditions that can be resolved through rollback
and retraining.
The four patterns of this section handle the conditions that
the Self-Healing Lifecycle cannot: the exhausted budget that
requires a regulatory decision, the stale circuit that
requires a GPU-day of recompilation, the broken composition
that requires a re-derivation of the stack guarantee, and
the decaying enclave that requires a coordinated
infrastructure operation.

Together, they form the complete operational framework for
a Trustworthy AI system that earns its trust not once,
at deployment, but continuously, in production.

\section{The Self-Healing Lifecycle}
\label{sec:self_healing}

A Trustworthy MLOps pipeline is not just a dashboard; it is an active control loop. We implement \textbf{Automated Fallbacks}:
\begin{enumerate}
    \item \textbf{Silent Monitoring:} The new model runs in parallel with a known-safe baseline.
    \item \textbf{Threshold Breach:} If drift or privacy risks exceed limits, the system diverts traffic back to the baseline model.
    \item \textbf{Automated Retraining:} The system initiates a retraining loop using curated, recently localized data (following the integrity rules of Chapter 24).
\end{enumerate}

% ------------------------------------------------------------------
\section{AI Model Forensics: Post-Incident Analysis and Evidence Preservation}
\label{sec:model_forensics}
% ------------------------------------------------------------------

Continuous monitoring and drift detection (Section~\ref{sec:self_healing}) tell you \emph{when} something goes wrong. \textbf{AI Model Forensics} tells you \emph{what}, \emph{why}, and \emph{who is accountable}---the structured investigation discipline that begins the moment an Emergency Brake fires or a compliance alert is raised. Without forensics, post-incident reviews are reconstructions from incomplete memory; with them, they are verifiable replays of a cryptographically signed chain of evidence.

Forensics in AI systems faces two constraints that do not appear in traditional software debugging:

\begin{description}
  \item[Privacy of Bystanders:] The evidence needed to reconstruct an incident---model inputs, intermediate activations, agent memory---often contains data belonging to \emph{other} users who were not party to the failure. Freezing and disclosing this evidence would constitute a secondary privacy breach.
  \item[Model Weight Volatility:] Unlike a static binary, model weights change continuously under retraining. Without a forensic snapshot, the model state at the moment of failure is irrecoverable within hours.
\end{description}

\subsection{Incident Classification}
\label{subsec:incident_classification}

Before a forensic protocol can be applied, the incident must be classified. Table~\ref{tab:incident_classes} defines four severity classes, each with distinct forensic obligations.

\begin{table}[ht]
\centering
\caption{AI Incident Classification and Forensic Obligations}
\label{tab:incident_classes}
\begin{tabular}{p{1.5cm} p{3.8cm} p{3.5cm} p{3.0cm}}
\hline\noalign{\smallskip}
\textbf{Class} & \textbf{Trigger} & \textbf{Evidence Required} & \textbf{Retention Period} \\
\noalign{\smallskip}\svhline\noalign{\smallskip}
S1 & Harm to a data subject (physical, financial, reputational) & Full forensic snapshot + ZK-replay + DPO notification & 7 years (GDPR Art.~5) \\[4pt]
S2 & Constitutional Policy violation without confirmed harm & Partial snapshot (context diff) + audit log chain & 3 years \\[4pt]
S3 & Drift threshold breach with no policy violation & Drift metrics + retraining trigger log & 1 year \\[4pt]
S4 & Operational anomaly (latency spike, OOM) & System telemetry + alert log & 90 days \\
\noalign{\smallskip}\hline\noalign{\smallskip}
\end{tabular}
\end{table}

\subsection{The Forensic Capture Protocol}
\label{subsec:forensic_capture}

For S1 and S2 incidents, the following four-phase Forensic Capture Protocol executes automatically upon brake acquisition (complementing Phase~0 of the HiC Transition Controller, Chapter~\ref{chap:hitl}):

\begin{enumerate}
  \item \textbf{Weight Snapshot.} The exact model checkpoint ($\theta_t$) active at the moment of the incident is hashed (SHA-256) and written to immutable cold storage, tagged with the incident ID and timestamp. This is the forensic ``crime scene.''
  \item \textbf{Differential Evidence Capture.} Only the context window of the \emph{violating session} is preserved---not raw inputs from other concurrent sessions. Inputs from bystander sessions are replaced with their HMAC tokens (verifiable proxies that confirm the input existed without revealing its content), preserving bystander privacy while maintaining audit integrity.
  \item \textbf{ZK-Replay Certificate.} A Zero-Knowledge proof is generated, attesting that the sequence of inputs in the frozen context window produces the recorded output under the frozen model weights---without revealing either the inputs or the weights to any inspector who does not hold the decryption key.
  \item \textbf{Regulatory Notification Package.} For S1 incidents, the system auto-generates a GDPR Article~33 notification draft, pre-populated with: the incident timestamp, the affected subject count (derived from HMAC token count), the technical root cause (the violated Datalog predicate), and the containment action taken (the HiC mutation applied).
\end{enumerate}

\begin{lstlisting}[language=Python, style=pythonstyle,
  caption={Forensic Capture Controller: preserves differential evidence for incident investigation while protecting bystander privacy through HMAC tokenization of concurrent session inputs.},
  label={lst:forensic-capture}]
import hashlib, time, json
from dataclasses import dataclass, field
from typing import Dict, List, Optional

@dataclass
class ForensicSnapshot:
    incident_id:        str
    severity_class:     str          # S1, S2, S3, S4
    weight_hash:        str          # SHA-256 of model checkpoint
    context_diff:       dict         # violating session only
    bystander_tokens:   List[str]    # HMAC proxies, not raw inputs
    zk_replay_cert:     Optional[str] = None
    timestamp:          float = field(default_factory=time.time)

class ForensicCaptureController:
    """
    Executes the four-phase Forensic Capture Protocol.
    Designed to be called from HiCTransitionController.acquire_brake().
    """
    def __init__(self, cold_storage, zk_prover,
                 hmac_key: bytes, audit_logger):
        self.storage   = cold_storage
        self.prover    = zk_prover
        self.hmac_key  = hmac_key
        self.auditor   = audit_logger

    def _hash_weights(self, checkpoint_path: str) -> str:
        import pathlib
        data = pathlib.Path(checkpoint_path).read_bytes()
        return hashlib.sha256(data).hexdigest()

    def _tokenize_bystanders(self, concurrent_sessions: List[dict]) \
            -> List[str]:
        """Replace raw bystander inputs with HMAC tokens."""
        return [
            hashlib.blake2b(
                json.dumps(s, sort_keys=True).encode(),
                key=self.hmac_key, digest_size=32
            ).hexdigest()
            for s in concurrent_sessions
        ]

    def capture(self, incident_id: str, severity: str,
                checkpoint_path: str, violating_context: dict,
                concurrent_sessions: List[dict]) -> ForensicSnapshot:
        # Phase 1: Weight snapshot
        w_hash = self._hash_weights(checkpoint_path)
        self.storage.write_immutable(
            key=f"forensic/{incident_id}/weights",
            value={"hash": w_hash, "path": checkpoint_path},
        )

        # Phase 2: Differential evidence
        bystander_tokens = self._tokenize_bystanders(
            concurrent_sessions)

        # Phase 3: ZK-Replay certificate (async in production)
        zk_cert = self.prover.prove_replay(
            context=violating_context,
            weight_hash=w_hash,
        ) if severity in ("S1", "S2") else None

        snap = ForensicSnapshot(
            incident_id      = incident_id,
            severity_class   = severity,
            weight_hash      = w_hash,
            context_diff     = violating_context,
            bystander_tokens = bystander_tokens,
            zk_replay_cert   = zk_cert,
        )

        self.auditor.log(event="forensic_capture_complete",
                         incident_id=incident_id,
                         severity=severity,
                         weight_hash=w_hash)
        return snap
\end{lstlisting}

\begin{svgraybox}
\textbf{Design Heuristic:} Evidence preservation and bystander privacy are not in tension if the forensic architecture is designed correctly from the start. The key principle is \textbf{differential evidence}: capture only the minimal context required to reproduce the incident, and tokenize everything else. A forensic snapshot that contains one session's full context plus HMAC proxies for all concurrent sessions satisfies both the regulator's demand for reproducibility and the data subjects' right to confidentiality.
\end{svgraybox}

\section{Conclusion}
\label{sec:conclusion_mlops}

Trustworthy MLOps ensures that our systems are resilient to the passage of time. By closing the loop between deployment and monitoring, we move from ``One-Time Trust'' to ``Continuous Accountability.'' With these operational safeguards in place, we can now navigate the global regulatory environment with the confidence that our systems remain compliant in practice, not just in theory. In the next chapter, we provide a comparative guide to \textbf{Global Compliance}.
